\item If $\vec{a}$, $\vec{b}$ and $\vec{c}$ are the position vectors of the points $A(2, 3, -4)$, $B(3, -4, -5)$ and $C(3, 2, -3)$ respectively, then $|\vec{a} + \vec{b} + \vec{c}|$ is equal to
\begin{enumerate}[label=(\Alph*), noitemsep]
\item $\sqrt{113}$
\item $\sqrt{185}$
\item $\sqrt{203}$
\item $\sqrt{209}$
\end{enumerate}
\hfill (12,2022)

\textbf{Solution:} The position vectors are the columns of
\begin{equation}
\mathbf{M} = \begin{pmatrix} \mathbf{a} & \mathbf{b} & \mathbf{c} \end{pmatrix} = \begin{pmatrix} 2 & 3 & 3 \\ 3 & -4 & 2 \\ -4 & -5 & -3 \end{pmatrix}
\end{equation}
The entries of the Gram matrix \eqref{c12eq:gramsum} are $G_{ij} = \mathbf{a}_i^{\top}\mathbf{a}_j$,
\begin{equation}
\mathbf{G} = \mathbf{M}^{\top}\mathbf{M} = \begin{pmatrix} 29 & 14 & 24 \\ 14 & 50 & 16 \\ 24 & 16 & 22 \end{pmatrix}
\end{equation}
Using \eqref{c12eq:gramsumnorm},
\begin{equation}
\|\mathbf{a} + \mathbf{b} + \mathbf{c}\|^2 = \mathbf{1}^{\top}\mathbf{G}\,\mathbf{1} = 29 + 50 + 22 + 2(14 + 24 + 16) = 209
\end{equation}
\begin{equation}
\therefore\ |\vec{a} + \vec{b} + \vec{c}| = \sqrt{209}
\end{equation}
Hence, the correct option is (D). The following code verifies this.

\nopagebreak[4]
\begin{center}
\fbox{\parbox{0.9\textwidth}{\centering 12/codes/q1.py}}
\end{center}

\item The distance between the planes $4x - 4y + 2z + 5 = 0$ and $2x - 2y + z + 6 = 0$ is
\begin{enumerate}[label=(\Alph*), noitemsep]
\item $\dfrac{1}{6}$
\item $\dfrac{7}{6}$
\item $\dfrac{11}{6}$
\item $\dfrac{16}{6}$
\end{enumerate}
\hfill (12,2022)

\textbf{Solution:} In the form $\mathbf{n}^{\top}\mathbf{x} = c$, the planes are
\begin{equation}
\begin{pmatrix} 4 & -4 & 2 \end{pmatrix}\mathbf{x} = -5, \qquad
\begin{pmatrix} 2 & -2 & 1 \end{pmatrix}\mathbf{x} = -6
\end{equation}
The normals have norms $6$ and $3$. Dividing each equation by the norm of its normal gives the normalised equations \eqref{c12eq:normalised},
\begin{equation}
\frac{1}{3}\begin{pmatrix} 2 & -2 & 1 \end{pmatrix}\mathbf{x} = -\frac{5}{6}, \qquad
\frac{1}{3}\begin{pmatrix} 2 & -2 & 1 \end{pmatrix}\mathbf{x} = -2
\end{equation}
Both have the same unit normal $\hat{\mathbf{n}} = \dfrac{1}{3}\begin{pmatrix} 2 \\ -2 \\ 1 \end{pmatrix}$, so the planes are parallel and the distance between them is
\begin{equation}
d = \left|p_1 - p_2\right| = \left|-\frac{5}{6} + 2\right| = \frac{7}{6}
\end{equation}
Hence, the correct option is (B). The following code verifies this.

\nopagebreak[4]
\begin{center}
\fbox{\parbox{0.9\textwidth}{\centering 12/codes/q3.py}}
\end{center}

\item The probability of solving a specific question independently by $A$ and $B$ are $\dfrac{1}{3}$ and $\dfrac{1}{5}$ respectively. If both try to solve the question independently, the probability that the question is solved is
\begin{enumerate}[label=(\Alph*), noitemsep]
\item $\dfrac{7}{15}$
\item $\dfrac{8}{15}$
\item $\dfrac{2}{15}$
\item $\dfrac{14}{15}$
\end{enumerate}
\hfill (12,2022)

\textbf{Solution:} The pmfs of the number of questions solved by $A$ and by $B$ alone are
\begin{equation}
\mathbf{p}_A = \begin{pmatrix} \frac{2}{3} \\ \frac{1}{3} \end{pmatrix}, \quad
\mathbf{p}_B = \begin{pmatrix} \frac{4}{5} \\ \frac{1}{5} \end{pmatrix}
\end{equation}
Let $X$ be the number of persons who solve the question. By \eqref{c12eq:pmfsum},
\begin{equation}
\mathbf{p}_X = \begin{pmatrix} \frac{2}{3}\cdot\frac{4}{5} \\ \frac{2}{3}\cdot\frac{1}{5} + \frac{1}{3}\cdot\frac{4}{5} \\ \frac{1}{3}\cdot\frac{1}{5} \end{pmatrix} = \frac{1}{15}\begin{pmatrix} 8 \\ 6 \\ 1 \end{pmatrix}, \quad \mathbf{1}^{\top}\mathbf{p}_X = 1
\end{equation}
The question is solved when at least one of them solves it,
\begin{equation}
\Pr(X \ge 1) = 1 - p_X(0) = 1 - \frac{8}{15} = \frac{7}{15}
\end{equation}
Hence, the correct option is (A). The following code verifies this.

\nopagebreak[4]
\begin{center}
\fbox{\parbox{0.9\textwidth}{\centering 12/codes/q4.py}}
\end{center}

\item $\displaystyle\int_0^1 \tan\left(\sin^{-1}x\right)dx$ equals
\begin{enumerate}[label=(\alph*), noitemsep]
\item $2$
\item $0$
\item $-1$
\item $1$
\end{enumerate}
\hfill (12,2022)

\textbf{Solution:} In $x$, the integrand $\dfrac{x}{\sqrt{1-x^2}} \to \infty$ at $x=1$, so \eqref{c12eq:trapz} cannot be applied directly. From \eqref{c12eq:sinsub},
\begin{equation}
I = \int_0^{\pi/2} \tan\theta\cos\theta\,d\theta = \int_0^{\pi/2} \sin\theta\,d\theta
\end{equation}
Taking a small step $h = \dfrac{\pi}{200}$ ($n = 100$), the nodes are $\theta_i = ih$, and from \eqref{c12eq:trapz},
\begin{equation}
I \approx h\left[\frac{\sin 0 + \sin\frac{\pi}{2}}{2} + \sum_{i=1}^{99} \sin\frac{i\pi}{200}\right]
= h\,\mathbf{w}^{\top}\mathbf{f} \approx 0.99998
\end{equation}
Decreasing $h$ gives
\begin{equation}
\begin{array}{c|cccc}
h & \frac{\pi}{12} & \frac{\pi}{200} & \frac{\pi}{400} & \frac{\pi}{800} \\ \hline
I & 0.994282 & 0.999979 & 0.999995 & 0.999999
\end{array}
\end{equation}
so $I \to 1$ as $h \to 0$. Hence, the correct option is (d). The following code computes this.

\nopagebreak[4]
\begin{center}
\fbox{\parbox{0.9\textwidth}{\centering 12/codes/q9.py}}
\end{center}

\item The roots of the equation $\begin{vmatrix} x & 0 & 8 \\ 4 & 1 & 3 \\ 2 & 0 & x \end{vmatrix} = 0$ are
\begin{enumerate}[label=(\alph*), noitemsep]
\item $-4, 4$
\item $2, -4$
\item $2, 4$
\item $2, 8$
\end{enumerate}
\hfill (12,2022)

\textbf{Solution:} Expanding along the second column, which has a single non-zero entry, using \eqref{c12eq:det},
\begin{equation}
\begin{vmatrix} x & 0 & 8 \\ 4 & 1 & 3 \\ 2 & 0 & x \end{vmatrix}
= 1\cdot\begin{vmatrix} x & 8 \\ 2 & x \end{vmatrix} = x^2 - 16 = 0
\end{equation}
Let $f(x) = x^2 - 16$, so $f'(x) = 2x$. From \eqref{c12eq:newton},
\begin{equation}
x_{n+1} = x_n - \frac{x_n^2 - 16}{2x_n} = \frac{1}{2}\left(x_n + \frac{16}{x_n}\right)
\end{equation}
Starting from $x_0 = 6$ and $x_0 = -6$,
\begin{equation}
\begin{array}{c|ccccc}
n & 0 & 1 & 2 & 3 & 4 \\ \hline
x_n & 6 & 4.3333 & 4.0128 & 4.00002 & 4.0000 \\
x_n & -6 & -4.3333 & -4.0128 & -4.00002 & -4.0000
\end{array}
\end{equation}
so the roots are $x = \pm 4$. Hence, the correct option is (a). The following code computes this and generates Fig. \ref{c12fig:q10}.

\begin{center}
\fbox{\parbox{0.9\textwidth}{\centering 12/codes/q10.py}}
\end{center}
\begin{figure}[!h]
\centering
\includegraphics[width=0.5\textwidth]{q10_newton.png}
\caption{Newton-Raphson iterates for $f(x)=x^2-16$ from $x_0=\pm6$}
\label{c12fig:q10}
\end{figure}

\item A card is picked at random from a pack of 52 playing cards. Given that the picked up card is a queen, the probability of it being a queen of spades is \underline{\hspace{1cm}}. \hfill (12,2022)

\textbf{Solution:} Let $X$ be the card picked. All 52 cards are equally likely, so
\begin{equation}
p_X(x) = \frac{1}{52}
\end{equation}
Let $A$ be the set of the 4 queens. Then
\begin{equation}
\Pr(X \in A) = \sum_{x \in A} p_X(x) = \frac{4}{52} = \frac{1}{13}
\end{equation}
From \eqref{c12eq:condpmf}, the pmf of $X$ given that it is a queen is $p_{X|A}(x) = \dfrac{1/52}{1/13} = \dfrac{1}{4}$ for each queen, so
\begin{equation}
p_{X|A}(\text{queen of spades}) = \frac{1}{4}
\end{equation}
Hence, the answer is $\dfrac{1}{4}$. The following code computes this.

\begin{center}
\fbox{\parbox{0.9\textwidth}{\centering 12/codes/q13.py}}
\end{center}

\item The integrating factor of the differential equation $x\dfrac{dy}{dx} - y = \log x$ is \underline{\hspace{1cm}}. \hfill (12,2022)

\textbf{Solution:} Here $\log$ is the natural logarithm. Writing $t$ for $x$ and taking the Laplace transform, using \eqref{c12eq:laplace},
\begin{equation}
-\frac{d}{ds}\left(sY\right) - Y = -\frac{\gamma + \ln s}{s}
\implies \frac{d}{ds}\left(s^2 Y\right) = \gamma + \ln s
\end{equation}
Integrating,
\begin{equation}
Y(s) = \frac{\gamma + \ln s - 1}{s} + \frac{C}{s^2}
\implies y(t) = Ct - \ln t - 1
\end{equation}
The $\dfrac{C}{s^2}$ term is the solution of the homogeneous equation, so $y_h = x$ for $C=1$. From \eqref{c12eq:ifactor},
\begin{equation}
\text{IF} = \frac{1}{y_h} = \frac{1}{x}
\end{equation}
Hence, the integrating factor is $\dfrac{1}{x}$. With $y(1) = 0$, $C = 1$ and
\begin{equation}
y = x - \ln x - 1
\end{equation}
Euler's forward difference method \eqref{c12eq:eulerfwd} is used to check this. With $f(x,y) = \dfrac{y + \ln x}{x}$, $x_0 = 1$, $y_0 = 0$, $h = 0.001$,
\begin{equation}
y_{n+1} = y_n + 0.001\,\frac{y_n + \ln x_n}{x_n}
\end{equation}
The two solutions are compared in Fig. \ref{c12fig:q14}, and at a few points
\begin{equation}
\begin{array}{c|ccc}
x & 1.0 & 1.5 & 2.0 \\ \hline
\text{Euler} & 0 & 0.0943 & 0.3065 \\
\text{Laplace} & 0 & 0.0945 & 0.3069
\end{array}
\end{equation}
For this small $h$, the Euler values are very close to the Laplace solution (the difference is below $4\times10^{-4}$). The following code computes both and generates Fig. \ref{c12fig:q14}.

\nopagebreak[4]
\begin{center}
\fbox{\parbox{0.9\textwidth}{\centering 12/codes/q14.py}}
\end{center}
\begin{figure}[!h]
\centering
\includegraphics[width=0.5\textwidth]{q14_stem.png}
\caption{Stem plot: Laplace solution and Euler forward difference ($h=0.001$)}
\label{c12fig:q14}
\end{figure}

\item Find the solution of the differential equation $\log\left(\dfrac{dy}{dx}\right) = ax + by$. \hfill (12,2022)

\textbf{Solution:} Here $\log$ is the natural logarithm, so
\begin{equation}
\frac{dy}{dx} = e^{ax + by}
\end{equation}
This is not linear in $y$. Let $u = e^{-by}$. Then
\begin{equation}
\frac{du}{dx} = -b\,e^{-by}\frac{dy}{dx} = -b\,e^{ax}
\end{equation}
Taking the Laplace transform using \eqref{c12eq:laplace2} and writing $u_0 = u(0)$,
\begin{equation}
sU - u_0 = -\frac{b}{s-a}
\implies U(s) = \frac{u_0}{s} - \frac{b}{s(s-a)} = \frac{u_0 + b/a}{s} - \frac{b/a}{s-a}
\end{equation}
using partial fractions. Taking the inverse transform,
\begin{equation}
u(x) = u_0 + \frac{b}{a} - \frac{b}{a}e^{ax} \implies e^{-by} = u_0 + \frac{b}{a} - \frac{b}{a}e^{ax}
\end{equation}
Dividing by $b$ and rearranging, the solution is
\begin{equation}
\frac{1}{a}e^{ax} + \frac{1}{b}e^{-by} = C
\end{equation}
To check this with Euler's forward difference method \eqref{c12eq:eulerfwd}, take $a = 2$, $b = -1$, $y(0) = 0$, so $u_0 = 1$ and the Laplace solution is $y = \ln\dfrac{e^{2x}+1}{2}$. With $h = 0.001$,
\begin{equation}
y_{n+1} = y_n + 0.001\,e^{2x_n - y_n}, \quad
\begin{array}{c|ccc}
x & 0 & 0.5 & 1.0 \\ \hline
\text{Euler} & 0 & 0.6199 & 1.4336 \\
\text{Laplace} & 0 & 0.6201 & 1.4338
\end{array}
\end{equation}
They agree closely, as shown in Fig. \ref{c12fig:q17}. The following code computes both and generates Fig. \ref{c12fig:q17}.

\nopagebreak[4]
\begin{center}
\fbox{\parbox{0.9\textwidth}{\centering 12/codes/q17.py}}
\end{center}
\begin{figure}[!h]
\centering
\includegraphics[width=0.5\textwidth]{q17_stem.png}
\caption{Stem plot: Laplace solution and Euler forward difference ($h=0.001$)}
\label{c12fig:q17}
\end{figure}


\item \textbf{(a)} If $\vec{a} = 2\hat{i} - \hat{j} + 2\hat{k}$ and $\vec{b} = 5\hat{i} - 3\hat{j} - 4\hat{k}$, then find the ratio of the projection of $\vec{a}$ on $\vec{b}$ to the projection of $\vec{b}$ on $\vec{a}$.

\textbf{(b)} Let $\hat{a}$ and $\hat{b}$ be unit vectors. If $\vec{c} = \hat{a} + 2\hat{b}$ and $\vec{d} = 5\hat{a} - 4\hat{b}$ are perpendicular to each other, then find the angle between $\hat{a}$ and $\hat{b}$. \hfill (12,2022)

\textbf{Solution (a):} In column vectors,
\begin{equation}
\mathbf{a} = \begin{pmatrix} 2 \\ -1 \\ 2 \end{pmatrix}, \quad
\mathbf{b} = \begin{pmatrix} 5 \\ -3 \\ -4 \end{pmatrix}
\implies \mathbf{a}^{\top}\mathbf{b} = 10 + 3 - 8 = 5, \quad
\|\mathbf{a}\| = 3, \quad \|\mathbf{b}\| = 5\sqrt{2}
\end{equation}
Using \eqref{c12eq:proj}, the required ratio is
\begin{equation}
\frac{\mathbf{a}^{\top}\mathbf{b}/\|\mathbf{b}\|}{\mathbf{b}^{\top}\mathbf{a}/\|\mathbf{a}\|}
= \frac{\|\mathbf{a}\|}{\|\mathbf{b}\|} = \frac{3}{5\sqrt{2}} = \frac{3\sqrt{2}}{10}
\end{equation}

The following code verifies this.

\nopagebreak[4]
\begin{center}
\fbox{\parbox{0.9\textwidth}{\centering 12/codes/q18a.py}}
\end{center}

\textbf{Solution (b):} Let $\mathbf{M} = \begin{pmatrix} \hat{\mathbf{a}} & \hat{\mathbf{b}} \end{pmatrix}$. Since $\|\hat{\mathbf{a}}\| = \|\hat{\mathbf{b}}\| = 1$, the Gram matrix is
\begin{equation}
\mathbf{G} = \mathbf{M}^{\top}\mathbf{M} = \begin{pmatrix} 1 & k \\ k & 1 \end{pmatrix}, \quad k = \hat{\mathbf{a}}^{\top}\hat{\mathbf{b}} = \cos\theta
\end{equation}
Now $\mathbf{c} = \mathbf{M}\begin{pmatrix} 1 \\ 2 \end{pmatrix}$ and $\mathbf{d} = \mathbf{M}\begin{pmatrix} 5 \\ -4 \end{pmatrix}$. By \eqref{c12eq:angle}, $\mathbf{c}\perp\mathbf{d}$ gives
\begin{equation}
\mathbf{c}^{\top}\mathbf{d} = \begin{pmatrix} 1 & 2 \end{pmatrix}
\begin{pmatrix} 1 & k \\ k & 1 \end{pmatrix}
\begin{pmatrix} 5 \\ -4 \end{pmatrix}
= 6k - 3 = 0 \implies \cos\theta = \frac{1}{2}
\end{equation}
\begin{equation}
\therefore\ \theta = \frac{\pi}{3}
\end{equation}
The following code verifies this.

\nopagebreak[4]
\begin{center}
\fbox{\parbox{0.9\textwidth}{\centering 12/codes/q18b.py}}
\end{center}


\item Find the distance of the point $(a, b, c)$ from the $x$-axis. \hfill (12,2022)

\textbf{Solution:} The $x$-axis passes through the origin with direction vector $\mathbf{m} = \mathbf{e}_1$. The foot of the perpendicular from $\mathbf{P}$ is the component of $\mathbf{P}$ along $\mathbf{m}$, so by \eqref{c12eq:distline}
\begin{equation}
\mathbf{P} = \begin{pmatrix} a \\ b \\ c \end{pmatrix}, \quad
\mathbf{Q} = \mathbf{m}\mathbf{m}^{\top}\mathbf{P} = \begin{pmatrix} a \\ 0 \\ 0 \end{pmatrix}
\end{equation}
The distance is the length of $\mathbf{P}-\mathbf{Q}$, which can be computed directly,
\begin{equation}
d = \|\mathbf{P}-\mathbf{Q}\| = \left\| \begin{pmatrix} 0 \\ b \\ c \end{pmatrix} \right\|
\end{equation}
\begin{equation}
\therefore\ d = \sqrt{b^2 + c^2}
\end{equation}
For example, for $\mathbf{P} = \begin{pmatrix} 2 \\ 3 \\ 4 \end{pmatrix}$, $\mathbf{Q} = \begin{pmatrix} 2 \\ 0 \\ 0 \end{pmatrix}$ and $d = 5$. The following code verifies this.

\nopagebreak[4]
\begin{center}
\fbox{\parbox{0.9\textwidth}{\centering 12/codes/q20.py}}
\end{center}


\clearpage
\item Show that the three vectors $2\hat{i} - \hat{j} + \hat{k}$, $\hat{i} - 3\hat{j} - 5\hat{k}$ and $3\hat{i} - 4\hat{j} - 4\hat{k}$ form the vertices of a right-angled triangle. \hfill (12,2022)

\textbf{Solution:} Let the position vectors of the vertices be
\begin{equation}
\mathbf{A} = \begin{pmatrix} 2 \\ -1 \\ 1 \end{pmatrix}, \quad
\mathbf{B} = \begin{pmatrix} 1 \\ -3 \\ -5 \end{pmatrix}, \quad
\mathbf{C} = \begin{pmatrix} 3 \\ -4 \\ -4 \end{pmatrix}
\end{equation}
The side vectors are the columns of
\begin{equation}
\mathbf{D} = \begin{pmatrix} \mathbf{B}-\mathbf{A} & \mathbf{C}-\mathbf{B} & \mathbf{A}-\mathbf{C} \end{pmatrix}
= \begin{pmatrix} -1 & 2 & -1 \\ -2 & -1 & 3 \\ -6 & 1 & 5 \end{pmatrix}
\end{equation}
Using \eqref{c12eq:gram},
\begin{equation}
\mathbf{G} = \mathbf{D}^{\top}\mathbf{D} = \begin{pmatrix} 41 & -6 & -35 \\ -6 & 6 & 0 \\ -35 & 0 & 35 \end{pmatrix}
\end{equation}
Since $\mathbf{G}_{23} = (\mathbf{C}-\mathbf{B})^{\top}(\mathbf{A}-\mathbf{C}) = 0$, the sides $BC$ and $CA$ are perpendicular. Hence $\triangle ABC$ is right-angled at $\mathbf{C}$. This also agrees with $\|\mathbf{B}-\mathbf{A}\|^2 = 41 = 6 + 35 = \|\mathbf{C}-\mathbf{B}\|^2 + \|\mathbf{A}-\mathbf{C}\|^2$, as shown in Fig. \ref{c12fig:q22}. The following code computes $\mathbf{G}$ and generates Fig. \ref{c12fig:q22}.

\nopagebreak[4]
\begin{center}
\fbox{\parbox{0.9\textwidth}{\centering 12/codes/q22.py}}
\end{center}
\begin{figure}[!h]
\centering
\includegraphics[width=0.4\textwidth]{q22_tri.png}
\caption{Triangle $ABC$ drawn in its own plane, right angle at $C$}
\label{c12fig:q22}
\end{figure}


\item \textbf{(a)} A bag contains 19 tickets, numbered 1 to 19. A ticket is drawn at random and then another ticket is drawn without replacing the first one in the bag. Find the probability distribution of the number of even numbers on the ticket.

\textbf{(b)} Find the probability distribution of the number of successes in two tosses of a die, when a success is defined as ``number greater than 5''. \hfill (12,2022)

\textbf{Solution (a):} Among the 19 tickets, 9 are even and 10 are odd. Let $X$ be the number of even tickets in the two draws, so $X \in \{0,1,2\}$. Let $E_i$ be the event that the $i$-th ticket is even. Using \eqref{c12eq:mult},
\begin{align}
p_X(0) &= \Pr(\bar{E}_1)\Pr(\bar{E}_2 \mid \bar{E}_1) = \frac{10}{19}\cdot\frac{9}{18} = \frac{5}{19} \\
p_X(1) &= \Pr(E_1)\Pr(\bar{E}_2 \mid E_1) + \Pr(\bar{E}_1)\Pr(E_2 \mid \bar{E}_1) = \frac{9}{19}\cdot\frac{10}{18} + \frac{10}{19}\cdot\frac{9}{18} = \frac{10}{19} \\
p_X(2) &= \Pr(E_1)\Pr(E_2 \mid E_1) = \frac{9}{19}\cdot\frac{8}{18} = \frac{4}{19}
\end{align}
In matrix form, with $\mathbf{p}_X = \begin{pmatrix} p_X(0) & p_X(1) & p_X(2) \end{pmatrix}^{\top}$,
\begin{equation}
\mathbf{p}_X = \frac{1}{19}\begin{pmatrix} 5 \\ 10 \\ 4 \end{pmatrix}, \quad \mathbf{1}^{\top}\mathbf{p}_X = 1
\end{equation}

The following code computes the pmf by listing all outcomes.

\nopagebreak[4]
\begin{center}
\fbox{\parbox{0.9\textwidth}{\centering 12/codes/q23a.py}}
\end{center}

\textbf{Solution (b):} A success is getting a $6$, so on each toss $\Pr(\text{success}) = \dfrac{1}{6}$. Let $Y$ be the number of successes in the two (independent) tosses, so $Y \in \{0,1,2\}$:
\begin{align}
p_Y(0) &= \frac{5}{6}\cdot\frac{5}{6} = \frac{25}{36}, \quad
p_Y(1) = 2\cdot\frac{1}{6}\cdot\frac{5}{6} = \frac{10}{36}, \quad
p_Y(2) = \frac{1}{6}\cdot\frac{1}{6} = \frac{1}{36} \\
\mathbf{p}_Y &= \frac{1}{36}\begin{pmatrix} 25 \\ 10 \\ 1 \end{pmatrix}, \quad \mathbf{1}^{\top}\mathbf{p}_Y = 1
\end{align}
The following code computes the pmf by listing all outcomes.

\nopagebreak[4]
\begin{center}
\fbox{\parbox{0.9\textwidth}{\centering 12/codes/q23b.py}}
\end{center}


\clearpage
\item Solve the differential equation: $x\dfrac{dy}{dx} = x + y$. \hfill (12,2022)

\textbf{Solution:} Rewriting, $x\dfrac{dy}{dx} - y = x$. Every solution found below has $y(0) = 0$. Let $y(x) \overset{\mathcal{L}}{\longleftrightarrow} Y(s)$. Applying $t\,g(t) \overset{\mathcal{L}}{\longleftrightarrow} -\dfrac{dG}{ds}$ from \eqref{c12eq:laplace} to $g = y'$ with \eqref{c12eq:laplace2},
\begin{equation}
x\,y' \overset{\mathcal{L}}{\longleftrightarrow} -\frac{d}{ds}\left(sY\right) = -sY' - Y
\end{equation}
Taking the Laplace transform of both sides, with $x \overset{\mathcal{L}}{\longleftrightarrow} \dfrac{1}{s^2}$,
\begin{equation}
-sY' - Y - Y = \frac{1}{s^2} \implies Y' + \frac{2}{s}Y = -\frac{1}{s^3}
\end{equation}
The homogeneous solution is $Y_h = s^{-2}$, so by \eqref{c12eq:ifactor} the integrating factor is $s^2$ and
\begin{equation}
\frac{d}{ds}\left(s^2 Y\right) = -\frac{1}{s} \implies Y(s) = \frac{C - \ln s}{s^2}
\end{equation}
Applying $t\,g(t) \overset{\mathcal{L}}{\longleftrightarrow} -\dfrac{dG}{ds}$ to $\ln x \overset{\mathcal{L}}{\longleftrightarrow} -\dfrac{\gamma+\ln s}{s}$ from \eqref{c12eq:laplace} gives $x\ln x \overset{\mathcal{L}}{\longleftrightarrow} \dfrac{1-\gamma-\ln s}{s^2}$, so
\begin{equation}
Y(s) = \frac{1-\gamma-\ln s}{s^2} + \frac{C-1+\gamma}{s^2}
\end{equation}
Taking the inverse transform, with $C_1 = C - 1 + \gamma$,
\begin{equation}
y = x\ln x + C_1 x
\end{equation}
To check this with Euler's forward difference method \eqref{c12eq:eulerfwd}, take $y(1) = 1$, so $C_1 = 1$ and the Laplace solution is $y = x\ln x + x$. With $h = 0.001$,
\begin{equation}
y_{n+1} = y_n + 0.001\,\frac{x_n + y_n}{x_n}, \quad
\begin{array}{c|ccc}
x & 1 & 1.5 & 2.0 \\ \hline
\text{Euler} & 1 & 2.1079 & 3.3858 \\
\text{Laplace} & 1 & 2.1082 & 3.3863
\end{array}
\end{equation}
They agree closely, as shown in Fig. \ref{c12fig:q26}. The following code computes both and generates Fig. \ref{c12fig:q26}.

\nopagebreak[4]
\begin{center}
\fbox{\parbox{0.9\textwidth}{\centering 12/codes/q26.py}}
\end{center}
\begin{figure}[!h]
\centering
\includegraphics[width=0.5\textwidth]{q26_stem.png}
\caption{Stem plot: Laplace solution and Euler forward difference ($h=0.001$)}
\label{c12fig:q26}
\end{figure}


\item The random variable $X$ has a probability function $P(x)$ as defined below, where $k$ is some number:
\begin{equation*}
P(x) = \begin{cases} k, & \text{if } x = 0 \\ 2k, & \text{if } x = 1 \\ 3k, & \text{if } x = 2 \\ 0, & \text{otherwise.} \end{cases}
\end{equation*}
Find: (i) the value of $k$; (ii) $P(X < 2)$, $P(X \le 2)$, $P(X \ge 2)$. \hfill (12,2022)

\textbf{Solution:} The pmf of $X \in \{0,1,2\}$ is the vector
\begin{equation}
\mathbf{p} = \begin{pmatrix} P(0) \\ P(1) \\ P(2) \end{pmatrix} = k\begin{pmatrix} 1 \\ 2 \\ 3 \end{pmatrix}
\end{equation}
(i) Since the probabilities sum to $1$,
\begin{equation}
\mathbf{1}^{\top}\mathbf{p} = 6k = 1 \implies k = \frac{1}{6}, \quad
\mathbf{p} = \frac{1}{6}\begin{pmatrix} 1 \\ 2 \\ 3 \end{pmatrix}
\end{equation}
(ii) Using \eqref{c12eq:cdf},
\begin{equation}
\begin{pmatrix} \Pr(X \le 0) \\ \Pr(X \le 1) \\ \Pr(X \le 2) \end{pmatrix}
= \begin{pmatrix} 1 & 0 & 0 \\ 1 & 1 & 0 \\ 1 & 1 & 1 \end{pmatrix}\frac{1}{6}\begin{pmatrix} 1 \\ 2 \\ 3 \end{pmatrix}
= \frac{1}{6}\begin{pmatrix} 1 \\ 3 \\ 6 \end{pmatrix}
\end{equation}
Therefore
\begin{equation}
P(X < 2) = P(X \le 1) = \frac{1}{2}, \quad
P(X \le 2) = 1, \quad
P(X \ge 2) = 1 - P(X < 2) = \frac{1}{2}
\end{equation}
The following code verifies these values.

\nopagebreak[4]
\begin{center}
\fbox{\parbox{0.9\textwidth}{\centering 12/codes/q28.py}}
\end{center}


\clearpage
\item \textbf{(a)} Find the shortest distance between the following lines and hence write whether the lines are intersecting or not.
\begin{equation*}
\frac{x-1}{2} = \frac{y+1}{3} = z, \qquad \frac{x+1}{5} = \frac{y-2}{1}, \ z = 2
\end{equation*}

\textbf{(b)} Find the equation of the plane through the line of intersection of the planes $\vec{r}\cdot(\hat{i} + 3\hat{j}) + 6 = 0$ and $\vec{r}\cdot(3\hat{i} - \hat{j} - 4\hat{k}) = 0$, which is at a unit distance from the origin. \hfill (12,2022)

\textbf{Solution (a):} Writing each line as $\mathbf{x} = \mathbf{a} + \lambda\mathbf{m}$ (the second line has $z = 2$, so its direction has zero $z$-component),
\begin{equation}
\mathbf{a}_1 = \begin{pmatrix} 1 \\ -1 \\ 0 \end{pmatrix}, \
\mathbf{m}_1 = \begin{pmatrix} 2 \\ 3 \\ 1 \end{pmatrix}, \quad
\mathbf{a}_2 = \begin{pmatrix} -1 \\ 2 \\ 2 \end{pmatrix}, \
\mathbf{m}_2 = \begin{pmatrix} 5 \\ 1 \\ 0 \end{pmatrix}
\end{equation}
Using \eqref{c12eq:skew},
\begin{equation}
\mathbf{n} = \mathbf{m}_1 \times \mathbf{m}_2 = \begin{vmatrix} \mathbf{e}_1 & \mathbf{e}_2 & \mathbf{e}_3 \\ 2 & 3 & 1 \\ 5 & 1 & 0 \end{vmatrix} = \begin{pmatrix} -1 \\ 5 \\ -13 \end{pmatrix}, \quad
\|\mathbf{n}\| = \sqrt{195}
\end{equation}
\begin{equation}
\mathbf{n}^{\top}(\mathbf{a}_2 - \mathbf{a}_1) = \begin{pmatrix} -1 & 5 & -13 \end{pmatrix}\begin{pmatrix} -2 \\ 3 \\ 2 \end{pmatrix} = 2 + 15 - 26 = -9
\end{equation}
\begin{equation}
\therefore\ d = \frac{|-9|}{\sqrt{195}} = \frac{9}{\sqrt{195}} \approx 0.64
\end{equation}
Since $d \neq 0$, the lines are \emph{not intersecting} (they are skew).

The following code verifies this.

\nopagebreak[4]
\begin{center}
\fbox{\parbox{0.9\textwidth}{\centering 12/codes/q30a.py}}
\end{center}

\textbf{Solution (b):} In the form $\mathbf{n}^{\top}\mathbf{x} = c$, the two planes are
\begin{equation}
\begin{pmatrix} 1 & 3 & 0 \\ 3 & -1 & -4 \end{pmatrix}\mathbf{x} = \begin{pmatrix} -6 \\ 0 \end{pmatrix}
\end{equation}
\emph{Step 1: line of intersection.} Using Gauss elimination \eqref{c12eq:gauss}, $R_2 \leftarrow R_2 - 3R_1$ gives
\begin{equation}
\left(\begin{array}{ccc|c} 1 & 3 & 0 & -6 \\ 3 & -1 & -4 & 0 \end{array}\right) \longrightarrow
\left(\begin{array}{ccc|c} 1 & 3 & 0 & -6 \\ 0 & -10 & -4 & 18 \end{array}\right)
\end{equation}
Taking the free variable $z = 5t$ and back-substituting, $y = -\frac{9}{5} - 2t$ and $x = -\frac{3}{5} + 6t$, so
\begin{equation}
\mathbf{x} = \mathbf{A} + t\,\mathbf{m}, \quad
\mathbf{A} = \begin{pmatrix} -\frac{3}{5} \\ -\frac{9}{5} \\ 0 \end{pmatrix}, \quad
\mathbf{m} = \begin{pmatrix} 6 \\ -2 \\ 5 \end{pmatrix}
\end{equation}
\emph{Step 2: the plane.} Let the required plane be $\mathbf{u}^{\top}\mathbf{x} = 1$ with $\|\mathbf{u}\| = 1$, so that its distance from the origin is $1$. It contains the line, so $\mathbf{u}^{\top}\mathbf{m} = 0$ and $\mathbf{u}^{\top}\mathbf{A} = 1$. With $\mathbf{u} = \begin{pmatrix} p & q & r \end{pmatrix}^{\top}$, and multiplying the second equation by $-5$,
\begin{equation}
\left(\begin{array}{ccc|c} 6 & -2 & 5 & 0 \\ 3 & 9 & 0 & -5 \end{array}\right) \xrightarrow{R_1 \leftarrow R_1 - 2R_2}
\left(\begin{array}{ccc|c} 0 & -20 & 5 & 10 \\ 3 & 9 & 0 & -5 \end{array}\right)
\end{equation}
Back-substituting with $q$ free,
\begin{equation}
r = 4q + 2, \quad p = -\frac{5}{3} - 3q
\end{equation}
Now $\|\mathbf{u}\| = 1$ gives
\begin{equation}
\left(\frac{5}{3} + 3q\right)^2 + q^2 + (4q+2)^2 = 1 \implies 9q^2 + 9q + 2 = 0 \implies q = -\frac{1}{3},\ -\frac{2}{3}
\end{equation}
\begin{equation}
q = -\tfrac{1}{3}:\ \mathbf{u} = \frac{1}{3}\begin{pmatrix} -2 \\ -1 \\ 2 \end{pmatrix}, \qquad
q = -\tfrac{2}{3}:\ \mathbf{u} = \frac{1}{3}\begin{pmatrix} 1 \\ -2 \\ -2 \end{pmatrix}
\end{equation}
Multiplying $\mathbf{u}^{\top}\mathbf{x} = 1$ by $3$, the planes are
\begin{equation}
\begin{pmatrix} 2 & 1 & -2 \end{pmatrix}\mathbf{x} + 3 = 0, \qquad
\begin{pmatrix} 1 & -2 & -2 \end{pmatrix}\mathbf{x} - 3 = 0
\end{equation}
that is, $\vec{r}\cdot(2\hat{i} + \hat{j} - 2\hat{k}) + 3 = 0$ and $\vec{r}\cdot(\hat{i} - 2\hat{j} - 2\hat{k}) - 3 = 0$. The following code verifies this.

\nopagebreak[4]
\begin{center}
\fbox{\parbox{0.9\textwidth}{\centering 12/codes/q30b.py}}
\end{center}


\clearpage
\item Find the equation of the line joining $A(1, 3)$ and $B(0, 0)$, using determinants. Also, find $k$ if $D(k, 0)$ is a point such that the area of $\triangle ABD$ is $3$ square units. \hfill (12,2022)

\textbf{Solution:} The points are
\begin{equation}
\mathbf{A} = \begin{pmatrix} 1 \\ 3 \end{pmatrix}, \quad
\mathbf{B} = \begin{pmatrix} 0 \\ 0 \end{pmatrix}, \quad
\mathbf{D} = \begin{pmatrix} k \\ 0 \end{pmatrix}
\end{equation}
A point $\mathbf{x} = \begin{pmatrix} x \\ y \end{pmatrix}$ lies on line $AB$ if $\mathbf{x}$, $\mathbf{A}$, $\mathbf{B}$ are collinear. By \eqref{c12eq:area}, expanding \eqref{c12eq:det} along the first row,
\begin{equation}
\begin{vmatrix} x & y & 1 \\ 1 & 3 & 1 \\ 0 & 0 & 1 \end{vmatrix} = 0
\implies x\begin{vmatrix} 3 & 1 \\ 0 & 1 \end{vmatrix} - y\begin{vmatrix} 1 & 1 \\ 0 & 1 \end{vmatrix} + \begin{vmatrix} 1 & 3 \\ 0 & 0 \end{vmatrix} = 0
\end{equation}
Writing the result as $\mathbf{n}^{\top}\mathbf{x} = c$, the equation of line $AB$ is
\begin{equation}
\begin{pmatrix} 3 & -1 \end{pmatrix}\mathbf{x} = 0
\end{equation}
For the area of $\triangle ABD$, using \eqref{c12eq:area},
\begin{equation}
\begin{vmatrix} 1 & 3 & 1 \\ 0 & 0 & 1 \\ k & 0 & 1 \end{vmatrix}
= 1(0 - 0) - 3(0 - k) + 1(0 - 0) = 3k
\end{equation}
\begin{equation}
\frac{1}{2}|3k| = 3 \implies |k| = 2 \implies k = \pm 2
\end{equation}
The following code verifies these values.

\nopagebreak[4]
\begin{center}
\fbox{\parbox{0.9\textwidth}{\centering 12/codes/q32.py}}
\end{center}


\item Evaluate $\displaystyle\int_1^3 \left(x^2 + 1 + e^x\right)dx$ as the limit of sums. \hfill (12,2022)

\textbf{Solution:} Here $f(x) = x^2 + 1 + e^x$, $a = 1$, $b = 3$. Dividing $[1, 3]$ into $n$ equal intervals of width $h = \dfrac{2}{n}$, the nodes are $x_i = 1 + ih$. By \eqref{c12eq:trapz}, the sum is
\begin{equation}
I_n = h\,\mathbf{w}^{\top}\mathbf{f}, \quad
\mathbf{w} = \begin{pmatrix} \frac{1}{2} & 1 & \cdots & 1 & \frac{1}{2} \end{pmatrix}^{\top}, \quad
\mathbf{f} = \begin{pmatrix} f(x_0) \\ \vdots \\ f(x_n) \end{pmatrix}
\end{equation}
and the integral is the limit $I = \lim_{n\to\infty} I_n$. Taking $h = 0.02$ ($n = 100$),
\begin{equation}
I \approx h\left[\frac{f(1) + f(3)}{2} + \sum_{i=1}^{99} f(1 + ih)\right] \approx 28.0346
\end{equation}
Decreasing $h$ gives
\begin{equation}
\begin{array}{c|cccc}
n & 4 & 10 & 100 & 1000 \\ \hline
h & 0.5 & 0.2 & 0.02 & 0.002 \\
I_n & 28.4776 & 28.1051 & 28.0346 & 28.0339
\end{array}
\end{equation}
so $I_n$ converges as $h \to 0$. The limit is
\begin{equation}
I = \left[\frac{x^3}{3} + x + e^x\right]_1^3 = \left(9 + 3 + e^3\right) - \left(\frac{1}{3} + 1 + e\right) = \frac{32}{3} + e^3 - e \approx 28.0339
\end{equation}
The following code computes $I_n$ for each $n$ and compares it with this value.

\nopagebreak[4]
\begin{center}
\fbox{\parbox{0.9\textwidth}{\centering 12/codes/q34.py}}
\end{center}
